%% Exercice : manipulation d'événements (TD1 Probabilités PGE1/S6 2025-2026, section 1)
%% Source de vérité pour l'énoncé, le corrigé et (plus tard) indices/pistes.

\ExoMeta{
  slug=proba-td1-evenements,
  titre=1 - Manipulation d'événements,
  td=proba-td1,
  niveau=1
}

\begin{Enonce}
On contrôle successivement une suite de pièces utilisées sur un chantier. Pour $n \geq 1$, on note $A_n$ l'événement \og la $n$-ième pièce contrôlée est conforme \fg.
\begin{enumerate}
    \item Comment écrire à partir des $A_i$ les événements :
    \begin{itemize}
        \item Les 3 premières pièces contrôlées sont conformes ;
        \item Parmi les 2 premières pièces contrôlées, la première est conforme et la seconde ne l'est pas ;
        \item Parmi les 2 premières pièces contrôlées, 1 est conforme et 1 ne l'est pas.
    \end{itemize}

    \item Décrire par une phrase les événements suivants :
    $$
    B=\bigcap_{i=5}^{+\infty} A_i, \qquad  C=\bigcup_{i=5}^{+\infty} A_i.
    $$

    \item Écrire à l'aide des $A_i$ l'événement $D$ \og À partir d'un certain contrôle, toutes les pièces contrôlées sont conformes \fg.

    \item On suppose que chaque pièce a une probabilité $\frac{1}{2}$ d'être conforme, indépendamment des autres. Calculer la probabilité que, parmi les deux premières pièces contrôlées, l'une au moins soit conforme.

    \item Quelle est la probabilité que la deuxième pièce soit conforme sachant que la première l'est ?

    \item Sachant qu'une des deux premières pièces est conforme, quelle est la probabilité que l'autre le soit aussi ?
\end{enumerate}
\end{Enonce}

\begin{Corrige}
\begin{enumerate}
    \item
    \begin{itemize}
        \item $A_1 \cap A_2 \cap A_3$
        \item $A_1 \cap \bar{A_2}$
        \item $(A_1 \cap \bar{A_2}) \cup (\bar{A_1} \cap A_2)$
    \end{itemize}

    \item $B$ signifie qu'à partir du cinquième contrôle, toutes les pièces contrôlées sont conformes. $C$ signifie qu'au moins une pièce contrôlée à partir du cinquième contrôle est conforme.

    \item Si on avait voulu l'événement $D_{10}$ : \og à partir du 10\ieme{} contrôle, toutes les pièces sont conformes \fg, on écrirait
    $$
    D_{10}=\bigcap_{i=10}^{+\infty} A_i.
    $$
    Pour écrire $D$, il faut considérer que c'est :
    \begin{enumerate}
        \item soit à partir du premier, toutes sont conformes,
        \item soit à partir du deuxième, toutes sont conformes,
        \item etc.
        \item soit à partir du dixième, toutes sont conformes,
        \item etc.
    \end{enumerate}
    C'est donc la réunion de tous ces événements :
    $$
    D=\bigcup_{n=1}^{+\infty} \bigcap_{i \geq n} A_i\,.
    $$

    \item On note $A_1$ l'événement \og la première pièce est conforme \fg{} et $A_2$ l'événement \og la deuxième pièce est conforme \fg. On cherche
    $$
    P(A_1 \cup A_2) = 1 - P(\bar{A_1} \cap \bar{A_2}).
    $$
    Par indépendance,
    $$
    P(\bar{A_1} \cap \bar{A_2}) = P(\bar{A_1}) \, P(\bar{A_2}) = \frac{1}{2} \times \frac{1}{2} = \frac{1}{4}.
    $$
    Donc
    $$
    P(A_1 \cup A_2) = 1 - \frac{1}{4} = \frac{3}{4}.
    $$

    \item Comme les contrôles sont indépendants,
    $$
    P(A_2 \mid A_1) = P(A_2) = \frac{1}{2}.
    $$
    Autrement dit, savoir que la première pièce est conforme ne change pas la probabilité que la deuxième le soit.

    \item L'information \og une des deux premières pièces est conforme \fg{} correspond à l'événement $A_1 \cup A_2$. Dire que l'autre l'est aussi signifie alors que les deux sont conformes, c'est-à-dire l'événement $A_1 \cap A_2$.

    On cherche donc
    $$
    P(A_1 \cap A_2 \mid A_1 \cup A_2) = \frac{P(A_1 \cap A_2)}{P(A_1 \cup A_2)}.
    $$
    Or, par indépendance,
    $$
    P(A_1 \cap A_2) = \frac{1}{2} \times \frac{1}{2} = \frac{1}{4},
    $$
    et, d'après la question précédente, $P(A_1 \cup A_2) = \frac{3}{4}$. Ainsi,
    $$
    P(A_1 \cap A_2 \mid A_1 \cup A_2) = \frac{\,\frac{1}{4}\,}{\,\frac{3}{4}\,} = \frac{1}{3}.
    $$
\end{enumerate}
\end{Corrige}

%% Les blocs suivants ne sont pas rendus dans le PDF : ils sont réservés à
%% l'extraction par le script de publication vers la SPA / l'IA.

\begin{Indices}
\begin{enumerate}
  \item Niveau 1 : pensez aux opérations ensemblistes classiques (intersection = \og et \fg, réunion = \og ou \fg, complémentaire = \og non \fg).
  \item Niveau 2 : pour \og exactement une pièce parmi deux \fg, décomposez en deux cas disjoints selon laquelle est conforme.
  \item Niveau 3 : pour l'événement $D$, commencez par écrire $D_n$ = \og toutes conformes à partir du rang $n$ \fg, puis prenez la réunion sur $n$.
\end{enumerate}
\end{Indices}

\begin{PlusLoin}
\begin{itemize}
  \item \textbf{Événement limite supérieur / inférieur} : les formules de $D$ et de son complémentaire font apparaître les notions de $\limsup A_n$ et $\liminf A_n$, utiles pour le lemme de Borel-Cantelli.
  \item \textbf{Indépendance mutuelle vs deux à deux} : dans la suite du TD on verra que l'indépendance mutuelle est strictement plus forte que l'indépendance deux à deux.
  \item \textbf{Schéma de Bernoulli} : cet exercice préfigure l'étude de la loi binomiale, qui compte le nombre de succès parmi $n$ épreuves indépendantes de même probabilité $p$.
\end{itemize}
\end{PlusLoin}
